\documentclass[11pt,a4paper]{article}
\usepackage[utf8]{inputenc}
\usepackage[T1]{fontenc}
\usepackage[portuguese]{babel}
\usepackage{xcolor}
\usepackage{hyperref}
\usepackage{geometry}
\usepackage{tcolorbox}

\geometry{margin=2.5cm}

\definecolor{primary}{RGB}{41, 128, 185}
\definecolor{secondary}{RGB}{44, 62, 80}
\definecolor{codebg}{RGB}{240, 240, 240}

\hypersetup{colorlinks=true, linkcolor=primary, urlcolor=primary}

\title{\color{secondary}\Huge\textbf{Exercício 5.4: Wikipedia com init e sidecar}}
\author{\color{primary}DevOps com Kubernetes -- 2026}
\date{}

\begin{document}
\maketitle

\section*{\color{primary}Instruções do Exercício}
\begin{tcolorbox}[colback=blue!5!white,colframe=blue!75!black,title=Exercício 5.4: Wikipedia com init e sidecar]
Escreva um aplicativo que sirva o conteúdo de páginas da Wikipedia a partir de \textit{containers} rodando no Kubernetes. A arquitetura do \textit{pod} deve conter:

- O contêiner principal baseado na imagem \textit{nginx}, configurado simplesmente para servir qualquer conteúdo HTML presente no seu diretório público padrão (\texttt{www}).

- Um contêiner de inicialização (\textit{init container}) que, no início do \textit{Pod}, execute um \texttt{curl} na página da Wikipedia sobre o Kubernetes e salve o conteúdo no diretório público montado no volume.

- Um \textit{sidecar container} rodando em paralelo que aguarde um tempo aleatório entre 5 e 15 minutos, execute um \texttt{curl} em uma página aleatória da Wikipedia (usando a URL de \textit{Special:Random}) e sobrescreva o conteúdo atual do diretório do nginx.
\end{tcolorbox}

\end{document}
